%%% Local Variables:
%%% mode: latex
%%% TeX-master: t
%%% End:

% ============================================================
% 第四章：图片与表格
% 本章介绍图片插入、表格制作和浮动体的使用方法。
% ============================================================

\chapter{图片与表格}\label{chap:floats}

\section{插入图片}

\subsection{基本用法}

图片使用 \texttt{figure} 浮动环境，\LaTeX{} 会自动选择最合适的位置放置：

\begin{verbatim}
\begin{figure}[htbp]
    \centering
    \includegraphics[width=0.6\textwidth]{你的图片文件名}
    \bicaption{中文图题}{English caption}
    \label{fig:示例}
\end{figure}
\end{verbatim}

实际效果如图~\ref{fig:ion}~所示：

\begin{figure}[htbp]
    \centering
    \includegraphics[width=0.6\textwidth]{ion}
    \bicaption{离子示意图}{Schematic diagram of ions}
    \label{fig:ion}
\end{figure}

如图~\ref{fig:lcn}~所示：

\begin{figure}[htbp]
    \centering
    \includegraphics[width=0.6\textwidth]{lcn}
    \bicaption{lcn示意图}{Schematic diagram of lcn}
    \label{fig:lcn}
\end{figure}

\textbf{参数说明：}

\begin{compactitem}
    \item \verb|[htbp]|：位置选项，依次尝试当前位置（h）、页面顶部（t）、底部（b）、单独页面（p）。
    \item \verb|\centering|：图片居中。
    \item \verb|width=0.6\textwidth|：图片宽度为版心宽度的60\%，可按需调整。
    \item \verb|\bicaption{中文}{英文}|：本模板要求中英双语图题（格式规范要求）。
    \item \verb|\label{fig:示例}|：标签，用于 \verb|\ref{fig:示例}| 引用。
\end{compactitem}

\textbf{图片格式建议：}
\begin{compactitem}
    \item 矢量图（折线图、柱状图等）：优先使用 \texttt{.pdf} 或 \texttt{.eps}，缩放不失真。
    \item 位图（照片、实验截图）：使用 \texttt{.png}（无损）或 \texttt{.jpg}，分辨率建议 $\geq$ 300 DPI。
    \item 文件名不含中文和空格，如 \texttt{detector\_setup.pdf}。
    \item 所有图片放在 \texttt{figures/} 目录下，引用时只写文件名，不写路径和扩展名。
\end{compactitem}

\subsection{调整图片大小}

\begin{verbatim}
\includegraphics[width=0.8\textwidth]{图片名}    % 按版心宽度比例
\includegraphics[height=5cm]{图片名}             % 按固定高度
\includegraphics[scale=0.5]{图片名}              % 按缩放比例
\end{verbatim}

\subsection{图的引用规范}

正文中必须有对图的提示，引用方式：

\begin{verbatim}
如图~\ref{fig:detector}~所示，探测器由三层构成。
图~\ref{fig:spectrum}~为实验测量的能谱。
\end{verbatim}

\section{三线表（推荐格式）}

格式规范推荐使用三线表：只有顶线、栏目线和底线，无竖线。
使用 \texttt{booktabs} 宏包提供的命令（已包含在模板中）：

\begin{verbatim}
\begin{table}[H]
\centering
\bicaption{表题中文}{Table caption in English}
\label{tab:示例}
\begin{tabular}{lccc}
\toprule                              % 顶线（粗）
列标题1 & 列标题2 & 列标题3 & 列标题4 \\
\midrule                              % 栏目线（细）
数据1   & 数据2   & 数据3   & 数据4   \\
数据5   & 数据6   & 数据7   & 数据8   \\
\bottomrule                           % 底线（粗）
\end{tabular}
\end{table}
\end{verbatim}

示例——实验参数表：

\begin{table}[H]
\centering
\bicaption{实验束流参数}{Beam parameters of the experiment}
\label{tab:beam}
\begin{tabular}{lcc}
\toprule
参数 & 数值 & 单位 \\
\midrule
束流能量 & 100 & MeV \\
束流强度 & $1 \times 10^6$ & pps \\
靶厚 & 2.0 & mg/cm$^2$ \\
探测角度 & 15--150 & 度 \\
\bottomrule
\end{tabular}
\end{table}

如表~\ref{tab:beam}~所示，束流能量为100 MeV。

\textbf{列格式说明（tabular参数）：}
\begin{compactitem}
    \item \texttt{l}：左对齐；\texttt{c}：居中；\texttt{r}：右对齐
    \item \texttt{p\{3cm\}}：固定宽度列，超出自动换行
    \item 竖线 \verb+|+ 可加列分隔线（三线表不使用）
\end{compactitem}

\section{带脚注的表格}

对表中数据需要补充说明时，使用 \texttt{threeparttable} 环境：

\begin{verbatim}
\begin{table}[H]
\centering
\bicaption{带脚注的表格示例}{Table with footnotes}
\begin{threeparttable}
\begin{tabular}{lcc}
\toprule
方法 & $\beta_3$ & $\delta_3$ (fm) \\
\midrule
电子散射\tnote{$a$} & 0.120 & 0.85 \\
$\alpha$散射\tnote{$b$}  & 0.118 & 0.84 \\
\bottomrule
\end{tabular}
\begin{tablenotes}
    \item[$a$] 束流能量75 MeV。
    \item[$b$] 束流能量79.1 MeV。
\end{tablenotes}
\end{threeparttable}
\end{table}
\end{verbatim}

\begin{table}[H]
\centering
\bicaption{$^{208}$Pb 形变参数的不同测量方法对比}{Comparison of $^{208}$Pb deformation parameters from different methods}
\begin{threeparttable}
\begin{tabular}{lcc}
\toprule
方法 & $\beta_3$ & $\delta_3$ (fm) \\
\midrule
电子散射\tnote{$a$} & 0.120 & 0.85 \\
$\alpha$散射\tnote{$b$} & 0.118 & 0.84 \\
\bottomrule
\end{tabular}
\begin{tablenotes}
    \item[$a$] 束流能量75 MeV。
    \item[$b$] 束流能量79.1 MeV。
\end{tablenotes}
\end{threeparttable}
\end{table}

\section{合并单元格}

使用 \verb|\multicolumn| 合并横向单元格，\verb|\multirow| 合并纵向单元格（需要 \texttt{multirow} 宏包）：

\begin{verbatim}
\multicolumn{列数}{对齐}{内容}    % 横向合并
\multirow{行数}{宽度}{内容}      % 纵向合并（需 \usepackage{multirow}）
\end{verbatim}

\section{图表的跨页处理}

当表格过长需要跨页，使用 \texttt{longtable} 环境：

\begin{verbatim}
\begin{longtable}{lcc}
\caption{跨页长表格}\label{tab:long} \\
\toprule
列1 & 列2 & 列3 \\
\midrule
\endfirsthead
% 续页表头（格式规范要求续表重复表头）
\multicolumn{3}{l}{\tablename~\thetable（续）} \\
\toprule
列1 & 列2 & 列3 \\
\midrule
\endhead
\bottomrule
\endfoot
% 表格内容行
数据1 & 数据2 & 数据3 \\
% ... 更多行
\end{longtable}
\end{verbatim}

\section{分图}

含有分图时，分图题置于各分图之下，分图号用 a)、b) 表示，整图标题置于所有分图之后。
使用 \texttt{subcaption} + \texttt{bicaption} 宏包（已在模板中载入），分图题支持中英双语：

\begin{verbatim}
\begin{figure}[htbp]
    \centering
    \begin{subfigure}{0.45\textwidth}
        \centering
        \includegraphics[width=\textwidth]{图1}
        \bisubcaption{中文分图题a}{English subcaption a}
        \label{fig:sub-a}
    \end{subfigure}
    \hfill
    \begin{subfigure}{0.45\textwidth}
        \centering
        \includegraphics[width=\textwidth]{图2}
        \bisubcaption{中文分图题b}{English subcaption b}
        \label{fig:sub-b}
    \end{subfigure}
    \bicaption{整图中文标题}{Overall English caption}
    \label{fig:subfig}
\end{figure}
\end{verbatim}

实际效果如图~\ref{fig:subfig-example}~所示：

\begin{figure}[htbp]
    \centering
    \begin{subfigure}{0.45\textwidth}
        \centering
        \includegraphics[width=\textwidth]{ion}
        \bisubcaption{离子示意图}{Schematic diagram of ions}
        \label{fig:sub-ion}
    \end{subfigure}
    \hfill
    \begin{subfigure}{0.45\textwidth}
        \centering
        \includegraphics[width=\textwidth]{lcn}
        \bisubcaption{lcn示意图}{Schematic diagram of lcn}
        \label{fig:sub-lcn}
    \end{subfigure}
    \bicaption{分图示例}{Example of subfigures}
    \label{fig:subfig-example}
\end{figure}

\section{本章小结}

本章介绍了图片插入、三线表制作、带脚注表格、合并单元格、
跨页表格以及分图的使用方法。
第~\ref{chap:references}~章将介绍参考文献的管理与引用。
